% LSPV-2 REPORT::FRONTBUF1::2026-10-12T07:39:11Z
% !TEX program = xelatex
\documentclass[11pt]{article}
\usepackage[T1]{fontenc}
\usepackage[margin=0.86in]{geometry}
\usepackage{amsmath,amssymb,amsthm,booktabs,longtable,array,enumitem,xurl,hyperref,microtype}
\newtheorem{theorem}{Theorem}
\newtheorem{corollary}{Corollary}
\newcommand{\SAT}{\mathsf{SAT}}
\newcommand{\EXACT}{\textbf{PROVED/EXACT}}
\newcommand{\KNOWN}{\textbf{KNOWN THEOREM}}
\newcommand{\COMP}{\textbf{COMPUTATIONAL}}
\newcommand{\COND}{\textbf{CONDITIONAL}}
\newcommand{\OPEN}{\textbf{OPEN}}
\newcommand{\RETIRED}{\textbf{FALSIFIED/RETIRED}}
\newcommand{\same}{\textsc{same\_realization}}
\title{FRONTBUF1: Explicit Clean-Core Buffers Miss the Calibrated FLMV Region}
\author{LSPV-2 / UNIVERSAL\_CORE\_TESTS}
\date{12 October 2026, 07:39 UTC}

\begin{document}
\maketitle

\begin{center}
\fbox{\parbox{0.94\linewidth}{
\textbf{Executive status.} P versus NP remains \OPEN.  KCALIB1's calibrated
composition formula is retained, but its illustrative zero-space front end
cannot be realized by first materializing and retaining a length-$M$ clean
core in ordinary writable memory.  Such a realization has both time and
space exponent at least the front-end length exponent $\lambda$, hence
$c+\rho\ge2\lambda$.  For the ordinary nonshrinking case $\lambda\ge1$ this
already exceeds the FLMV golden-ratio threshold.  Thus
\texttt{EXPLICIT-BUFFER-CALIBRATED-LANDING1} is \RETIRED.  Streaming,
virtual, and fused implementations are not ruled out, but each requires a
new charged generator/access theorem.}}
\end{center}

\section{Claim card, predecessor, and exact target}
\begin{description}[leftmargin=0cm,labelwidth=0cm]
\item[Claim ID.] \texttt{FRONTBUF1 / EXPLICIT-RETAINED-FRONTEND-FLOOR1}.
\item[Status before run.] The clean-core front-end access cost in KCALIB1 was
explicitly declared chargeable, but the illustrative ledger used $v=0$.
\item[Claim attacked.] A calibrated favorable child exponent by itself makes
an explicit retained clean-core realization compatible with the FLMV landing.
\item[Exact prefix.] For every ordinary same-machine composition satisfying
the explicit-retained-buffer hypotheses below and every finite solver bounds
$(K,D)$, its charged exponent pair satisfies $c+\rho\ge2\lambda$.
\item[Input/model.] An ordinary encoded CNF of length $N$ on a deterministic
random-access machine; writable working storage counts as space.  The front
end produces an encoded object of length $M=N^{\lambda+o(1)}$.
\item[Resource.] Total time and peak writable workspace of one \same{}
execution, including front end, retained buffer, compiler, calls, and output.
\item[What falsifies the theorem.] A realization satisfying the stated
explicit-retention hypothesis but writing or retaining $M$ independently
addressable bits in $o(M)$ time or writable space.
\end{description}

The immediately preceding LSPV-1 report is
\path{research/agents/lspv-1/2026-10-12_0721_LSPV-1_KCALIB1.tex}.  It was read
back from GitHub, is substantive, and contains both its required header and
final sentinel.  No predecessor write recovery is needed.

\section{Primary-literature and obstacle preflight}
The two obstacles are (i) realizing the front-end output inside the same
machine without an uncharged $M$-bit store and (ii) placing that same machine
in the simultaneous time--space region required by FLMV.  FLMV is a
\textbf{directly transferable theorem} only after the complete ordinary
machine has the stated random-access bounds~\cite{FLMV05}.  Hong--Kung
pebbling is a \textbf{useful structural analogy requiring a new transfer
theorem}: it separates retained storage from data movement in a declared
memory hierarchy, but is not an unrestricted SAT lower bound~\cite{HK81}.
VLSI, communication, proof-system, streaming, and cell-probe lower bounds are
\textbf{misleading analogies / failed transfers} without a model-preserving
compiler.  Relativization and algebrization remain active barriers
\cite{BGS75,AW09}.

\section{Exact retained-buffer floor}
KCALIB1 considers a front end $R$ with
\[
 |R(x)|=M=N^{\lambda+o(1)},\qquad
 T_R(N)=N^{u+o(1)},\qquad S_R(N)=N^{v+o(1)},
\]
followed by a clean-core compiler and calls to a fixed exact decider.  We now
type one concrete realization.

\begin{theorem}[\texttt{EXPLICIT-RETAINED-FRONTEND-FLOOR1}]
Suppose the same ordinary machine writes all $M$ output bits of $R(x)$ into
an explicitly addressed writable buffer and keeps those $M$ bits live until
the next phase may access them.  In a standard bit-cost model,
\[
 u\ge\lambda,qquad v\ge\lambda.
\]
Consequently every full composition has
\[
 c\ge\lambda,\qquad \rho\ge\lambda,qquad c+\rho\ge2\lambda. \tag{1}
\]
The conclusion is independent of the fixed solver exponents $(K,D)$ and the
child exponent $\beta$.
\end{theorem}
\begin{proof}
Writing $M$ explicit bits takes at least $M$ bit writes, so total time is
$\Omega(M)$.  At the instant the retained representation is complete and
still live, the writable buffer contains $M$ bits, so peak writable workspace
is $\Omega(M)$.  Substituting $M=N^{\lambda+o(1)}$ gives the exponent floors.
All later phases can only increase total time and peak space.
\end{proof}

\begin{corollary}[FLMV exclusion for the explicit-buffer realization]
For the FLMV target $c+\rho<\varphi=(1+\sqrt5)/2$, explicit retention requires
\[
 \lambda<\varphi/2\approx0.809. \tag{2}
\]
Hence every explicit retained front end with $\lambda\ge1$ misses the target
before the compiler or SAT calls are considered.
\end{corollary}

This does not prove that every reduction has $\lambda\ge1$, nor does it rule
out a sublinear kernel.  It says exactly that the usual linear-or-larger
explicit intermediate cannot instantiate KCALIB1's illustrative $v=0$
ledger.  A write-only output tape is not a counterexample: if the later phase
cannot read it, it is not the required accessible intermediate; if it can be
read as writable/random-access storage, its charge belongs to the model.

\section{Smallest repair: a realization trichotomy}
Any viable calibration must specify one of the following rather than writing
``produce $F$'' and assigning zero access space.
\begin{longtable}{p{0.19\linewidth}p{0.72\linewidth}}
\toprule Realization & Required charge or missing theorem\\ \midrule
Explicit retained buffer & Equation (1); misses FLMV for $\lambda\ge1$.\\
Streaming & Prove the compiler consumes the reduction stream in order with
small live state; charge every pass, stalled producer state, and output.\\
Virtual random access & Prove a bit generator; charge every requested bit,
address state, repeated generation, caching, parent retention, and
recomputation.\\
Fused reduction/compiler & Prove one source-to-children transducer directly;
charge recognition, normalization, failed routes, child access, and replay.\\
External/read-only representation & State the machine model and simulation;
do not import a free random-access store absent from the FLMV model.\\
\bottomrule
\end{longtable}

This is consistent with \path{main5.tex}, which already warns that regenerated
child access is not free.  FRONTBUF1 moves the same guard one level earlier:
the clean-core reduction itself cannot leave a free retained intermediate.

\section{Bounded reproducible audit}
The dependency-free script
\path{experiments/review/explicit_frontend_buffer_floor_audit.py} completed in
one sub-30-second chunk and wrote
\path{results/raw/review/2026-10-12_explicit_frontend_buffer_floor.json}.
It checked all $16{,}384$ tuples
\[
 \lambda\in\{1,11/10,5/4,3/2\},\qquad K,D\in\{1,\ldots,64\},
\]
and confirmed $c+\rho\ge2\lambda>\varphi$ with zero failures.  The finite
copy ledger used seven formula families and seven representations/wrappers:
49 rows, 29,152 charged bit writes, zero extraction/reconstruction calls, and
zero communication bits.  Controls included Horn, 2SAT, affine/Tseitin,
bounded treewidth, independent products, padding, redundant clauses, DIMACS,
JSON, COPY, DEAD, DELAY, INTERP, and RECOMP.  Each row separately records a
streaming line buffer, a virtual descriptor, and a fused zero-buffer control;
these are outside the theorem's hypothesis, not claimed landings.  Status is
\COMP{} only; no asymptotic inference is drawn from byte counts.

Script SHA-256:
\nolinkurl{80f3339d417df9b78cb5db0881447a97565ae433398d4f2acdc51d357005114e}.
Raw JSON SHA-256:
\nolinkurl{98c8ee710179926f4fb13e76354bafaa908e29d4bda62036aa5079eb900ed6ec}.

\section{Hostile-control and barrier audit}
Easy Horn, 2SAT, affine/Tseitin, and bounded-treewidth inputs can be consumed
by specialized streaming or algebraic algorithms; they do not refute the
explicit-buffer theorem and do not prove a universal compiler.  Independent
products, redundant clauses, padding, and alternate encodings enlarge the
materialized representation and therefore strengthen its finite buffer
charge, while also warning that length must be measured after encoding.
One-bit decision output forces no retained reduction object.  COPY and DELAY
cannot improve (1); INTERP and RECOMP may avoid storage only by moving work to
generation/access, which must be charged.

The theorem relativizes and algebrizes and yields no Natural-Proofs property.
Schur squares, quadratic relations, Hilbert functions, syzygies, and graded
Betti numbers are \textbf{FAILED TRANSFERS} because no code representation is
forced.  Collision entropy is an \textbf{EXACT OBJECT/THEOREM TRANSFER} only
for a fixed map and has no bearing on whether the intermediate is buffered.
LWE-style external extraction and special-soundness collision extraction are
\textbf{ANALOGY/CANDIDATE} routes requiring a new unconditional bridge.
Cryptographic binding and random-oracle assumptions are inadmissible.

\section{CORPUS REVIEWED}
The run surveyed the full repository tree; read root \path{main.tex},
\path{AGENTS.md}, \path{CODEX.md}, the full research skill and red-team/
experiment protocols, every task specification, and actual \path{main4.tex}
and \path{main5.tex}.  It verified the newest MAIN/SYNTHESIS-0022 and current
KCALIB1, COVCOLL1, MINFRONT1, COREDTCH1, RESOURCEAXIS1, HEVALLB1, APIFLAT1,
TESTS-0014--0021, and relevant earlier buffer/virtual-child reports.  It
inspected HEC/VSC-HEC/G-PPC-HEC/CRS-PPC/CCSPC/RCCN explorations, governance,
outboxes, and cited experiment artifacts.  Repository search covered
calibration, explicit children, virtual inputs, retained parents, streaming,
recomputation, padding, and solver-exponent dependence.  FLMV, Hong--Kung,
BGS, and Aaronson--Wigderson were checked in primary sources before promotion.

\section{DUPLICATION / CONTRADICTION AUDIT}
TESTS-0016 already proved that a small virtual child need not be cheaply
locatable; STREAMCOMP1 already charges a supplied bit generator.  KCALIB1
already says $v$ includes the actual output-access realization.  New here is
the exact simultaneous floor $c+\rho\ge2\lambda$ for the explicit retained
front-end buffer and its direct FLMV consequence (2).  There is no
contradiction: KCALIB1's algebra is retained, while one concrete realization
of its parameters is excluded.  Main5's explicit-child option concerns
strictly smaller children and does not grant a free full-scale front-end
buffer.

\section{WHAT CHANGED RELATIVE TO ROOT MAIN}
Root MAIN should retain \texttt{K-CALIBRATED-CLEANCORE-LANDING1} but qualify
the illustrative $v=0$ example: it requires a proved streaming, virtual, or
fused front end.  An explicit retained clean-core buffer with $\lambda\ge1$
cannot land below FLMV.  The overall status remains \OPEN.

\section{ADOPTED / REJECTED / RETIRED PRIOR COMPONENTS}
\textbf{Adopted:} KCALIB1's composition formula, encoded-length accounting,
STREAMCOMP1/VIRTUAL-CHILD-SPACE1 charges, FLMV model matching, and \same{}.

\textbf{Rejected:} free random access to a materialized reduction output;
counting only child size while omitting its acquisition/access; and treating a
write-only output tape as later random-access workspace.

\textbf{FALSIFIED/RETIRED:}
\texttt{EXPLICIT-BUFFER-CALIBRATED-LANDING1} for $\lambda\ge1$.  The broader
calibrated streaming/virtual/fused compiler remains \OPEN.

\section{STALE OR OVERSTATED CLAIMS FOUND}
Any use of KCALIB1's illustrative ledger with $v=0$ is overstated unless the
front end's access realization is proved.  ``The reduction is polynomial''
does not imply a favorable simultaneous time--space exponent.  No verified
construction or source theorem is downgraded.

\section{Status and handoff}
\begin{center}\small
\begin{tabular}{p{0.45\linewidth}p{0.18\linewidth}p{0.28\linewidth}}
\toprule Item & Status & Reason\\ \midrule
\texttt{EXPLICIT-RETAINED-FRONTEND-FLOOR1} & \EXACT & Write/live-buffer counting\\
FLMV endpoint & \KNOWN & Exact model required\\
Finite hostile-control audit & \COMP & 16,384 tuples; 49 rows\\
\texttt{EXPLICIT-BUFFER-CALIBRATED-LANDING1} & \RETIRED & $2\lambda>\varphi$ for $\lambda\ge1$\\
Streaming/virtual/fused calibrated compiler & \OPEN & No source theorem\\
Barrier escape & \OPEN & Result relativizes/algebrizes\\
P versus NP & \OPEN & Load-bearing compiler absent\\
\bottomrule
\end{tabular}
\end{center}

\paragraph{Precise missing implication.}
Construct one source-only clean-core-to-children transducer for each fixed
polynomial SAT decider and valid exponent bounds such that either (a) it
streams without retaining a full-scale intermediate, (b) it supplies a
charged favorable virtual bit generator, or (c) it fuses reduction and
compilation.  Prove exact all-input semantics and an FLMV landing after
charging recognition, generation, every access, reread, recomputation,
parent retention, calls, replay, output, time, and peak workspace in one
\same{} realization.  This remains separation-strength and \OPEN.

\section{Artifacts}
\begin{itemize}
\item \path{experiments/review/explicit_frontend_buffer_floor_audit.py}
\item \path{results/raw/review/2026-10-12_explicit_frontend_buffer_floor.json}
\item \path{research/review/outbox/tests/2026-10-12_0739_FRONTBUF1.md}
\item \path{research/review/logs/tests/2026-10-12_0739_FRONTBUF1.json}
\end{itemize}

\begin{thebibliography}{9}
\bibitem{FLMV05} L. Fortnow, R. Lipton, D. van Melkebeek and A. Viglas,
``Time-Space Lower Bounds for Satisfiability,'' \emph{J. ACM} 52(6),
835--865, 2005. \url{https://doi.org/10.1145/1101821.1101822}.
\bibitem{HK81} J.-W. Hong and H. T. Kung,
``I/O Complexity: The Red-Blue Pebble Game,'' \emph{STOC}, 1981,
326--333. \url{https://doi.org/10.1145/800076.802486}.
\bibitem{BGS75} T. Baker, J. Gill and R. Solovay,
``Relativizations of the $P=?NP$ Question,'' \emph{SIAM J. Comput.} 4(4),
431--442, 1975. \url{https://doi.org/10.1137/0204037}.
\bibitem{AW09} S. Aaronson and A. Wigderson, ``Algebrization: A New Barrier
in Complexity Theory,'' \emph{ACM TOCT} 1(1), 2009.
\url{https://eccc.weizmann.ac.il/report/2008/005/}.
\end{thebibliography}

% END_OF_VERIFIED_REPORT::LSPV-2::FRONTBUF1
\end{document}
